\documentclass[12pt]{article}
\usepackage[utf8]{inputenc}
\usepackage[T1]{fontenc}
\usepackage{amsmath,amssymb,amscd}
\usepackage{geometry}
\geometry{margin=2.5cm}

\begin{document}

\section*{Lecture 3 -- Page 4}

\noindent Let $(\mathbb{Z},+)$ be a commutative group, $(M,\cdot,e)$ a monoid.

\noindent We have $f : \mathbb{Z} \to M$, $f(k) = a^k$

\[
f(h+k) = f(h)\cdot f(k),
\]
i.e.\ $a^{h+k} = a^h \cdot a^k$.

\[
\begin{aligned}
h&=m, & k&=n\\
h&=-m, & k&=-n
\end{aligned}
\qquad m,n\in\mathbb{N}^*.
\]

\[
a^{-m}\cdot a^{-n} = (a^{-1})^{m}\cdot(a^{-1})^{n} = (a^{-1})^{m+n} = a^{h+k}
\]
\noindent\textit{(last term unclear in the original -- could be $a^{h+k}$ or $a^h\cdot a^k$; please check.)}

\bigskip
\noindent\underline{\textbf{Composition law induced on a stable subset}} (respectively on a finite set)

\bigskip
\noindent\textbf{Def.} Let $(S,\varphi)$; $\varphi : S\times S \to S$, $S \neq \varnothing$.

\noindent A subset $T \subseteq S$ is called a \textbf{stable subset} of $S$ with respect to $\varphi$ if $(\forall)\, x,y \in T \Rightarrow \varphi(x,y) \in T$.

\[
\begin{CD}
S\times S @>\varphi>> S\\
@AAA @AAA\\
T\times T @>\varphi'>> T
\end{CD}
\]
\noindent $(x,y)\in T \Rightarrow \varphi'(x,y) \overset{\text{def}}{=} \varphi(x,y)$

\noindent $\varphi'$ is the operation induced on $T$ by $\varphi$.

\noindent Let $S = M_2(\mathbb{R})$, $(M_2(\mathbb{R}),\cdot)$

\[
T = \left\{ \begin{pmatrix} 1 & a \\ 0 & 1 \end{pmatrix} \,\middle|\, a\in\mathbb{R} \right\} \subseteq S = M_2(\mathbb{R}).
\]
\[
X_a X_b = X_{a+b}: \qquad \begin{pmatrix}1 & a\\ 0 & 1\end{pmatrix}\begin{pmatrix}1 & b\\ 0 & 1\end{pmatrix} = \begin{pmatrix}1 & a+b\\ 0 & 1\end{pmatrix}
\]

\bigskip
\noindent\fbox{Proposition} \ Let $S\neq\varnothing$, $\varphi$ a composition law on $S$.\\
$T\subseteq S$, $T$ stable with respect to $\varphi$,\\
$\varphi'$ = the operation induced by $\varphi$ on $T$: $\varphi'(x,y)=\varphi(x,y)\in T$.

\begin{enumerate}
\item[(1)] If $\varphi$ is associative (resp.\ commutative) $\Rightarrow$ $\varphi'$ is associative (resp.\ commutative).
\item[(2)] If $\varphi$ admits an identity element $e$ and $e\in T$ $\Rightarrow$ $e$ is also an identity element for $\varphi'$.
\item[(3)] If an element $x\in T$ is invertible with respect to $\varphi$ ($\varphi$ associative with an identity element) and its inverse $x'$ with respect to $\varphi$ lies in $T$ $\Rightarrow$ $x$ is invertible with respect to $\varphi'$ as well.
\item[(4)] its inverse (with respect to $\varphi'$) is still $x'$.
\end{enumerate}

\end{document}
